\chapter{Real QDM data}\label{ch:real}

Everything so far was tested on simulation and on COMSOL. On 16 November Ritaban sent six maps measured on a real
quantum diamond microscope, of the same four-coil chip that the COMSOL model describes. This chapter shows what it takes
to use them, how every method did, and what can honestly be concluded. Nothing was trained and nothing was chosen on
these data.

\section{What a QDM actually measures}

The maps are 165 $\times$ 165 pixels of 1\um. Each pixel is \emph{one} number: the field projected on the axis of
the NV centres in the diamond,
\begin{equation}
B_{NV} = \mathbf{u}\cdot\mathbf{B} = u_x B_x + u_y B_y + u_z B_z ,
\end{equation}
where $\mathbf{u}$ is the NV axis. Every method in this booklet expects $B_x$, $B_y$, $B_z$. The NV axis, the standoff
and the currents were not given.

\begin{figure}[h]
\centering
\includegraphics[width=\textwidth]{rd_experimental_maps}
\caption{The six measured maps. Cases 1 and 4: bottom-left coil. Case 2: top-left coil (the mirror image, as in the chip
design). Case 5: both left coils. Cases 3 and 6: a coil seen from much farther away.}
\end{figure}

\section{A close look at the maps}

\begin{figure}[h]
\centering
\includegraphics[width=\textwidth]{rd_data_characterisation}
\caption{Top: local noise of case 1, the coil's signal spectrum against the noise floor, and how alike the cases are.
Bottom: case 1, the fitted design, and what the design does not explain.}
\end{figure}

\begin{itemize}
\item \textbf{Noise:} about 0.37\,\textmu T\ per pixel in every case, 3--8\% of the coil signal's range; the corners are
several times noisier, and a few hundred pixels per map are fit failures of up to $\pm10^5$\,\textmu T.
\item \textbf{Seam:} the left and right halves of each image differ by up to 1\,\textmu T\ at $x\approx86$\um.
\item \textbf{What the data contain:} the coil's spectrum meets the noise at about 0.85\,rad/\um\ (a 7\um\ wavelength)
for the near coils and 0.55\,rad/\um\ (11\um) for the deep ones. Finer detail is not in the data.
\item \textbf{Same coil, different cases:} cases 1, 4 and 5 correlate 0.91--0.96 over the bottom-left coil; case 3 is
the same shape from farther away.
\item \textbf{Leads:} the real leads run left along $y\approx20$--25\um, not where the drawing puts them.
\end{itemize}

\section{One projection is enough, if you know the axis}

Above the wires the field is source-free, so all three components come from one scalar potential, and in Fourier space
$B_x = -i(k_x/k)B_z$ and $B_y = -i(k_y/k)B_z$. Substituting,
\begin{equation}
B_z(\mathbf{k}) = \frac{B_{NV}(\mathbf{k})}{u_z - i\,(u_x k_x + u_y k_y)/k}.
\end{equation}
The denominator never falls below $u_z$, so the step is stable. But it needs $\mathbf{u}$.

\begin{figure}[h]
\centering
\includegraphics[width=\textwidth]{rd_vector_from_bnv_case1}
\caption{Case 1: $B_x$, $B_y$, $B_z$ rebuilt from the single projection (top) and the fitted design (bottom). This is
the input every method receives.}
\end{figure}

\section{Fitting the design to find the axis}

We placed the designed spiral (the COMSOL drawing) in each image and fitted its orientation, centre, standoff, the NV
axis, the current and a background. The fit finds the chip layout by itself: the top-left coil comes out mirrored, as
designed. Fitted standoffs: 3.4, 3.7, 4.3\um\ for the near coils and 9.2, 10.6\um\ for cases 3 and 6.

\begin{careful}[Two limits of this fit]
\begin{itemize}
\item The axis is weakly fixed by any one map: changing it moves the share of the signal explained only from 49 to
57\% (case 1). A joint scan over all maps at the standard angle $\theta = 54.7^\circ$ picks $\varphi = 157.5^\circ$.
\item The real leads are routed differently from the drawing, so the fitted design explains only 43--56\% of the near
maps. It is the reference we score against, and it is a \emph{model}, not ground truth.
\end{itemize}
\end{careful}

\section{The results}

\begin{figure}[h]
\centering
\includegraphics[width=\textwidth]{rd_magic_currents_all_cases}
\caption{MAGIC-UNet run 10b (0.7\um\ pixels, 44.8\um\ view) on all six maps. Bottom: predicted current. Green: the
designed spiral at the fitted position.}
\end{figure}

\begin{center}
\small
\begin{tabular}{lccccc}
\toprule
0.5\um\ field vs fitted design, in band & MAGIC-UNet & FISTA & FISTA-Net & CSC-LISTA & Tikhonov \\
\midrule
Case 1 (3.4\um) & 70\% & 76\% & 78\% & 69\% & \textbf{60\%} \\
Case 4 (4.3\um) & 80\% & 136\% & 111\% & 105\% & \textbf{73\%} \\
Case 3 (9.2\um) & 78\% & \textbf{49\%} & 88\% & 100\% & -- \\
\bottomrule
\end{tabular}
\end{center}

\begin{keyidea}[What the real data say]
\begin{itemize}
\item \textbf{The current layout comes through.} Run 10b draws the spiral in all six cases, including two coils at
once and coils at 9--12\um, without being told the standoff or the number of coils.
\item \textbf{The near-field advantage does not transfer.} Against the fitted design, Tikhonov is closest at
3--4\um\ and FISTA at 9\um. No learned method beats them on real data yet.
\item \textbf{The pixel must match the data.} Runs 2 and 3 (0.2\um\ pixels) turn the camera noise, which is smooth over
1\um\ pixels, into false currents and read the height far too low. Run 10b's 0.7\um\ pixels do not.
\item \textbf{The height comes from the prior.} Each network reads a height inside its own training range, as
Chapter~\ref{ch:question} predicted.
\end{itemize}
\end{keyidea}

\begin{figure}[h]
\centering
\includegraphics[width=\textwidth]{rd_methods_case1}
\caption{Case 1, $B_z$ at 0.5\um\ from every method, next to the fitted design.}
\end{figure}

\section{Is the pipeline right?}

On the clean field of the fitted coil, the same pipeline reads the standoff as 3.19\um\ (fit 3.37), reconstructs the
0.5\um\ field to 14\% and returns a current that correlates 0.985 with the coil. So the gap on real data is the data:
correlated noise, an uncertain axis, a background and leads that the simulator never showed the network.

\begin{defend}[If your professor asks: so does it work on real data?]
It recovers where the current flows, from 3 to 12\um, without being told the standoff; that part works. It does not
yet give a better near field than Tikhonov against our best reference, which is itself a fitted drawing. To measure
accuracy properly we need the NV axis, the standoff and the currents from the experiment, and a calibration structure
of known shape.
\end{defend}
